\newtheorem{papertheorem}{Theorem}
\newtheorem{paperlemma}{Lemma}[section]
\newtheorem{paperproposition}[paperlemma]{Proposition}
\newtheorem{papercorollary}[paperlemma]{Corollary}
\theoremstyle{definition}
\newtheorem{paperdefinition}[paperlemma]{Definition}
% Reject a mismatch between a statement number on GitHub and its printed number.
\newcommand{\checkstatement}[2]{%
  \begingroup\edef\actualnumber{#1}\def\sourcenumber{#2}%
  \ifx\actualnumber\sourcenumber\else
    \PackageError{manuscript}{Statement #2 has TeX number \actualnumber}
      {Check the order and section numbers of the Markdown theorem statements.}%
  \fi\endgroup}
% Read both fonts by filename from the TeX bundle, so code listings build
% independently of the host system's font configuration.
\font\leansource="[DejaVuSansMono.ttf]" at 8.2pt
\font\leanmath="[latinmodern-math.otf]" at 8.2pt
\newcommand{\leanchar}[1]{%
  \iffontchar\leansource"#1\relax\char"#1\relax
  \else\iffontchar\leanmath"#1\relax{\leanmath\char"#1\relax}%
  \else\PackageError{manuscript}{Missing Lean source glyph U+#1}{Supply a font with this glyph.}%
  \fi\fi}
\newcommand{\leangap}{\hskip\fontdimen2\leansource\relax}
% Keep a complete declaration on the page whenever it can fit on one page.
% The smaller strut supplies the leading of the code font instead of the prose font.
\newbox\leanlisting
\newenvironment{leanscript}{%
  \par\setbox\leanlisting=\vbox\bgroup\hsize=\linewidth
  \addvspace{5pt}\begingroup\leansource\baselineskip=10.5pt
  \setbox\strutbox=\hbox{\vrule height 7.3pt depth 2pt width 0pt}%
  \parskip=0pt\parindent=0pt\rightskip=0pt plus 1fil\parfillskip=0pt
  \hyphenpenalty=10000\exhyphenpenalty=10000
}{%
  \par\endgroup\addvspace{6pt}\egroup
  \ifdim\dimexpr\ht\leanlisting+\dp\leanlisting\relax>\textheight
    \unvbox\leanlisting
  \else\box\leanlisting\fi\par}
\newcommand{\leanline}[1]{\noindent\hangindent=2em\hangafter=1\strut#1\par}
\newcommand{\proofsources}[1]{{\normalfont\footnotesize\color{proofink}#1}}
\providecommand{\tightlist}{\setlength{\itemsep}{0pt}\setlength{\parskip}{0pt}}
\providecommand{\pandocbounded}[1]{#1}
